Software vulnerabilities continue to pose a significant threat to the security of modern systems,
with thousands of new Common Vulnerabilities and Exposures (CVEs) published each year.
Manual code review is costly and does not scale; automated static analysers produce high false-positive
rates and lack cross-language portability. This dissertation investigates whether a pre-trained
code model can be fine-tuned to detect vulnerabilities reliably across multiple programming languages,
and examines the gap between performance on clean synthetic data and performance on real-world code.

A multi-language synthetic dataset covering 17 CWE types across C and Python is constructed,
including eight shared cross-language vulnerability classes.
UniXcoder is fine-tuned as a binary classifier on this dataset, achieving F1\,=\,0.931 and
AUC-ROC\,=\,0.964 on a held-out synthetic validation set of 1{,}230 samples.
A cross-language ablation study reveals that joint training consistently outperforms
single-language training, and that transfer quality depends on the structural similarity
of a vulnerability's surface representation across languages rather than its conceptual equivalence.

A semantic understanding test suite probes whether the model is performing keyword matching
or genuine pattern recognition. The model is found to be robust to variable renaming and
keyword ablation, detecting most vulnerability patterns from structural context alone,
but it cannot reason about control-flow reachability or cross-statement data flow ---
an inherent limitation of sequence-level encoding.

Two targeted interventions address the model's principal failure modes.
Hard negative mining resolves systematic false positives on safe cryptographic functions
by adding contrastive training pairs.
Real-world data augmentation with 1{,}000 DiverseVul samples reduces the
synthetic-to-real F1 gap from $-0.522$ to $-0.182$, raising real-world performance
from 0.410 to 0.749.
The results establish a clear baseline, identify the architectural ceiling of sequence-based
approaches, and motivate graph-augmented architectures as the primary direction for future work.
